\setcounter{section}{27}

  \zwischenueberschrift{Turunan vertikal sepanjang suatu pemetaan}

Misalkan
\maabb {p} { E } { M
} {}
suatu
\definitionsverweis {bundel vektor diferensiabel}{}{}
di atas suatu
\definitionsverweis {manifold diferensiabel}{}{}
$M$ yang dilengkapi dengan suatu
\definitionsverweis {koneksi linear}{}{.}
Koneksi ini menghasilkan suatu
\definitionsverweis {turunan vertikal}{}{.}
Untuk setiap medan vektor
\maabb {V} { M } { TM
} {}
dan setiap penampang yang terdiferensialkan secara kontinu
\maabbdisp {s} { M } { E
} {}
\mathl{\nabla_V s}{} adalah penampang kontinu di $E$ yang diberikan oleh

\mathdisp {M \stackrel{V}{ \longrightarrow  } TM  \stackrel{T(s) }{ \longrightarrow } TE \stackrel{  \pi_{\text{vert} }  }{ \longrightarrow} E} {  }

Pada suatu himpunan terbuka

\mavergleichskette
{\vergleichskette
{ U
}
{ \subseteq }{ M
}
{  }{
}
{    }{
}
{   }{
}
}
{}{}{}
dengan

\mavergleichskette
{\vergleichskette
{ U
}
{ \cong }{ V
}
{ \subseteq }{ \R^n
}
{    }{
}
{   }{
}
}
{}{}{}
serta turunan-turunan parsial
\mathl{\partial_1  , \ldots , \partial_n}{} dan suatu trivialisasi

\mavergleichskettedisp
{\vergleichskette
{ E \vert_U
}
{ \cong} { U \times \R^r
}
{ } {
}
{ } {
}
{ } {
}
}
{}{}{}
dengan penampang basis
\mathl{s_1 , \ldots , s_r}{,} turunan vertikal semacam itu dideskripsikan sepenuhnya oleh simbol-simbol Christoffel $\Gamma_{ij}^k$ melalui

\mavergleichskettedisp
{\vergleichskette
{ \nabla_{\partial_i } (s_j)
}
{ =} { \sum_{k= 1}^r  \Gamma_{ij}^k s_k
}
{ } {
}
{ } {
}
{ } {
}
}
{}{}{}
berdasarkan
Lema 25.10.
Kita ingin memiliki konsep ini bukan hanya bagi penampang di atas $M$, melainkan juga bagi penampang
\maabbdisp {s} { L } { E
} {}
sepanjang suatu
\definitionsverweis {pemetaan diferensiabel}{}{}
tetap
\maabbdisp {\psi} { L } { M
} {}
yaitu ketika terdapat diagram komutatif

\mathdisp {\begin{matrix}  &   &   &  E   \\ & & s \nearrow & \downarrow  \\  &  L  & \stackrel{ \psi }{\longrightarrow}  &  M \! \\ \end{matrix}} {  }

Situasi ini telah kita jumpai pada penampang horizontal. Perhatikan bahwa suatu penampang
\maabbdisp {s} { L } { E
} {}
sama dengan suatu penampang
\maabbdisp {\tilde{s}} { L  } { \psi^*(E) = E \times_XY
} {}
di dalam
\definitionsverweis {bundel vektor tarik balik}{}{}
$\psi^*(E)$. Untuk suatu medan vektor $F$ pada $L$ dan suatu penampang semacam itu, kita dapat meninjau komposisi pemetaan

\mathdisp {L \stackrel{F}{ \longrightarrow  } TL  \stackrel{T(s) }{ \longrightarrow } TE \stackrel{  \pi_{\text{vert} }  }{ \longrightarrow} E} {  }

yang kembali kita nyatakan dengan $\nabla_F s$. Koneksi yang ditentukan oleh turunan vertikal ini kita sebut
\stichwort {koneksi tarik balik} {} $\psi^*\nabla$. Koneksi tersebut mempunyai sifat-sifat berikut.



\inputfaktbeweis
{Linearer Zusammenhang/Vertikale Ableitung/Längs Abbildung/Eigenschaften/Fakt}
{Lema}
{}
{

\faktsituation {Misalkan
\maabb {p} { E } { M
} {}
suatu
\definitionsverweis {bundel vektor diferensiabel}{}{}
di atas suatu
\definitionsverweis {manifold diferensiabel}{}{}
$M$ yang dilengkapi dengan suatu
\definitionsverweis {koneksi linear}{}{}
$\nabla$. Misalkan
\maabb {\psi} { L } { M
} {}
suatu
\definitionsverweis {pemetaan diferensiabel}{}{}
dengan
\definitionsverweis {koneksi tarik balik}{}{}
$\psi^*\nabla$ pada
\definitionsverweis {bundel vektor tarik balik}{}{}
$\psi^*E$ di atas $L$. Maka berlaku sifat-sifat berikut.}
\faktfolgerung {\aufzaehlungzwei {Pemetaan
\maabbeledisp {\psi^*\nabla} {C^1(\psi^* E)} { C^0(\psi^* E \otimes T^* L )
} { s } { \left(  \psi^*\nabla  \right) (s)
} {,}
bersifat
$\R$-\definitionsverweis {linear}{}{.}
Secara khusus, $\psi^*\nabla$ merupakan suatu koneksi linear.
} {Untuk suatu pembatasan
\maabbdisp {\psi} {U'} {U
} {}
pada domain-domain peta terbuka dengan koordinat
\mathl{z_1 , \ldots , z_m}{} di $U'$ dan koordinat
\mathl{x_1 , \ldots , x_n}{}
di $U$, bagi
\definitionsverweis {simbol-simbol Christoffel}{}{}
$\tilde{\Gamma}_{\ell j}^k$ dari $\psi^*\nabla$ dan
\definitionsverweis {penampang basis}{}{}

\mathl{s_1  , \ldots , s_r}{} berlaku hubungan

\mavergleichskettedisp
{\vergleichskette
{ \tilde{\Gamma}_{j \ell }^k
}
{ =} { \sum_{i = 1}^n  \left(  \partial_j \psi_i  \right)   \Gamma_{i\ell }^k
}
{ } {
}
{ } {
}
{ } {
}
}
{}{}{.}
}}
\faktzusatz {}
\faktzusatz {}

}
{

\aufzaehlungzwei {Jelas.
} {Kita tinjau situasi lokal secara langsung. Suatu
\definitionsverweis {penampang basis}{}{}
$s_\ell$ sepanjang $L$ dapat dipandang langsung sebagai penampang basis di atas $M$. Pemetaan-pemetaan yang relevan adalah

\mathdisp {U' \stackrel{\partial_j }{\longrightarrow} TU'  \stackrel{T (\psi) }{\longrightarrow} T(U) \stackrel{ T(s_\ell ) }{\longrightarrow} T( U \times \R^r) \stackrel{  \pi_{\text{vert} } }{\longrightarrow} \R^r} { . }

Di sini,

\mavergleichskettedisp
{\vergleichskette
{ T(\psi) \circ \partial_j
}
{ =} { \partial_j \psi
}
{ =} { \sum_{i = 1}^n \left(  \partial_j \psi_i  \right)  \partial_i
}
{ } {
}
{ } {
}
}
{}{}{.}
Dengan demikian, untuk

\mavergleichskette
{\vergleichskette
{ Q
}
{ \in }{ U'
}
{  }{
}
{    }{
}
{   }{
}
}
{}{}{}
dan dengan menggunakan
Lema 25.10,

\mavergleichskettealignhandlinks
{\vergleichskettealignhandlinks
{ \left(  \left(  \psi^*\nabla  \right) _{\partial_j } (s_\ell)  \right) (Q)
}
{ =} { \left(  \psi^*\nabla  \right) _{\partial_j (Q) } (s_\ell) (Q)
}
{ =} { \nabla_{  \left(  T(\psi) \circ \partial_j  \right)  (Q)  }  (s_\ell) ( \psi(Q) )
}
{ =} { \nabla_{  \sum_{i = 1}^n  \left(  \partial_j \psi_i  \right)  (Q) \partial_i  }  (s_\ell) ( \psi(Q) )
}
{ =} { \sum_{i = 1}^n  \left(  \partial_j \psi_i  \right)  (Q)\nabla_{  \partial_i }  (s_\ell) ( \psi(Q) )
}
}
{
\vergleichskettefortsetzungalign
{ =} { \sum_{i = 1}^n  \left(  \partial_j \psi_i  \right)  (Q)  \left(  \sum_{k = 1}^r \Gamma_{i \ell }^k(\psi(Q)) s_k (\psi(Q))  \right)
}
{ =} { \sum_{k = 1}^r  \left(  \sum_{i = 1}^n  \left(  \partial_j \psi_i  \right)  (Q)  \Gamma_{i \ell }^k(\psi(Q))   \right) s_k (\psi(Q))
}
{ } {}
{ } {}
}
{}{.}
Dengan meninjau komponen ke-$k$, kita memperoleh

\mavergleichskettedisp
{\vergleichskette
{ \tilde{\Gamma}_{j \ell }^k (Q)
}
{ =} { \sum_{i = 1}^n \Gamma_{i\ell }^k (\psi(Q))  \left(  \partial_j \psi_i  \right)  (Q)
}
{ } {
}
{ } {
}
{ } {
}
}
{}{}{.}
}

}



\inputfaktbeweis
{Linearer Zusammenhang/Vertikale Ableitung/Längs Abbildung/Metrisch/Fakt}
{Lema}
{}
{

\faktsituation {Misalkan
\maabb {p} { E } { M
} {}
suatu
\definitionsverweis {bundel vektor diferensiabel}{}{}
di atas suatu
\definitionsverweis {manifold diferensiabel}{}{}
$M$ yang dilengkapi dengan suatu
\definitionsverweis {koneksi linear}{}{}
\definitionsverweis {metrik}{}{}
$\nabla$. Misalkan
\maabb {\psi} { L } { M
} {}
suatu
\definitionsverweis {pemetaan diferensiabel}{}{}
dengan
\definitionsverweis {koneksi tarik balik}{}{}
$\psi^*\nabla$ pada
\definitionsverweis {bundel vektor tarik balik}{}{}
$\psi^*E$ di atas $L$.}
\faktfolgerung {Maka
\mathl{\psi^*\nabla}{} juga metrik.}
\faktzusatz {}
\faktzusatz {}

}
{

Pertama, menurut
Soal 16.13,
$\psi^*E$ kembali merupakan suatu bundel Riemann. Kita tinjau situasi lokal
\maabbdisp {\psi} {U'} {U
} {}
dan

\mavergleichskette
{\vergleichskette
{ E
}
{ = }{ U \times \R^r
}
{  }{
}
{    }{
}
{   }{
}
}
{}{}{}
serta, bersesuaian dengan itu,

\mavergleichskette
{\vergleichskette
{ \psi^*E
}
{ = }{ U' \times \R^r
}
{  }{
}
{    }{
}
{   }{
}
}
{}{}{.}
Fungsi-fungsi yang mendeskripsikan struktur Riemann pada $E$,

\mavergleichskettedisp
{\vergleichskette
{ h_{ \ell m}
}
{ =} { \left\langle s_\ell , s_m   \right\rangle
}
{ } {
}
{ } {
}
{ } {
}
}
{}{}{,}
pada $U$ dan $U'$ berhubungan langsung melalui $\psi$. Misalkan $\partial_j$ suatu
\definitionsverweis {medan vektor standar}{}{}
pada $U'$ dan
\mathl{s_\ell, s_m}{} penampang basis di $E$. Maka

\mavergleichskettealign
{\vergleichskettealign
{ \partial_j  \left\langle s_\ell , s_m   \right\rangle
}
{ =} { \partial_j  \left(  h_{\ell , m} \circ\psi   \right)
}
{ =} { \sum_{i = 1}^n  \left(  \partial_j \psi_i  \right) \left(  \partial_i  h_{\ell , m}  \right)
}
{ =} { \sum_{i = 1}^n  \left(  \partial_j \psi_i  \right) \partial_i  \left\langle s_\ell , s_m   \right\rangle
}
{ } {}
}
{}
{}{.}
Menurut
Lema 27.1,

\mavergleichskettealignhandlinks
{\vergleichskettealignhandlinks
{ \left\langle  \left(  \psi^*\nabla  \right)_{\partial_j } s_\ell ,  s_m   \right\rangle
}
{ =} { \left\langle  \sum_{k = 1}^r \tilde{\Gamma}_{j \ell }^k s_ k ,  s_m   \right\rangle
}
{ =} { \sum_{k = 1}^r \left\langle  \tilde{\Gamma}_{j \ell }^k s_ k ,  s_m   \right\rangle
}
{ =} { \sum_{k = 1}^r \left\langle  \sum_{i = 1}^n \Gamma_{i\ell }^k  \left(  \partial_j \psi_i  \right)  s_k  ,  s_m   \right\rangle
}
{ =} { \sum_{i = 1}^n \left(  \partial_j \psi_i  \right) \left\langle  \sum_{k = 1}^r \Gamma_{i\ell }^k s_k  ,  s_m   \right\rangle
}
}
{
\vergleichskettefortsetzungalign
{ =} { \sum_{i = 1}^n \left(  \partial_j \psi_i  \right) \left\langle  \nabla_{\partial_i } s_\ell  ,  s_m   \right\rangle
}
{ } {}
{ } {}
{ } {}
}
{}{}
dan, secara bersesuaian,

\mavergleichskettedisp
{\vergleichskette
{ \left\langle  s_\ell ,  \left(  \psi^*\nabla  \right)_{\partial_j } s_m   \right\rangle
}
{ =} { \sum_{i = 1}^n  \left(  \partial_j \psi_i  \right)   \left\langle  s_\ell  , \nabla_{\partial_i } s_m   \right\rangle
}
{ } {
}
{ } {
}
{ } {
}
}
{}{}{.}
Karena $\nabla$ metrik, diperoleh

\mavergleichskettedisp
{\vergleichskette
{ \partial_j  \left\langle s_\ell , s_m   \right\rangle
}
{ =} { \left\langle  \left(  \psi^*\nabla  \right)_{\partial_j } s_\ell ,  s_m   \right\rangle    + \left\langle  s_\ell ,  \left(  \psi^*\nabla  \right)_{\partial_j } s_m   \right\rangle
}
{ } {
}
{ } {
}
{ } {
}
}
{}{}{}
dan dari sini, bersama
Soal 26.10,
diperoleh bahwa $\psi^*\nabla$ metrik.

}

  \zwischenueberschrift{Percepatan tangensial}

Misalkan $M$ suatu
\definitionsverweis {manifold}{}{.}
Untuk suatu
\definitionsverweis {kurva diferensiabel}{}{}
\maabbdisp {\gamma} { I } { M
} {}
berlaku bahwa
\maabbdisp {\gamma'} { I} {TM
} {}
merupakan kurva di dalam bundel tangen yang menjadi suatu pengangkatan dari $\gamma$. Jika kurva ini pada gilirannya terdiferensialkan, diperoleh suatu kurva
\maabbdisp {\gamma^{\prime \prime}} { I} {TTM
} {,}
namun kurva ini tidak dapat langsung dihubungkan dengan $\gamma'$ karena nilainya berada di ruang lain. Sebaliknya, apabila pada $TM$ diberikan suatu
\definitionsverweis {koneksi}{}{},
maka melalui turunan vertikal
\zusatzklammer {tarik balik} {} {}
turunan kedua dapat didefinisikan sebagai
\mathl{\nabla_{ \partial } \gamma'}{}.

\inputdefinition
{}
{

Misalkan $M$ suatu
\definitionsverweis {manifold Riemann}{}{}
yang dilengkapi dengan
\definitionsverweis {koneksi Levi--Civita}{}{.}
Untuk suatu
\definitionsverweis {kurva diferensiabel}{}{}
\maabbdisp {\gamma} { I } { M
} {}
yang dua kali terdiferensialkan secara kontinu, pemetaan
\mathl{\pi_{\text{vert} }  \circ TT(\gamma)}{} disebut
\definitionswort {percepatan tangensial}{}
dari $\gamma$.

}

Besaran ini juga ditulis sebagai
\mathl{\pi_{\text{vert} } \circ \gamma^{\prime \prime}}{} atau
\zusatzklammer {sebagai turunan vertikal sepanjang $\gamma$} {} {}

\mathl{\nabla_\partial \gamma'}{,} dengan
\maabbdisp {\gamma'} { I } { TM
} {}
dipandang sebagai penampang dalam bundel tangen sepanjang lintasan $\gamma$, dan dengan meninjau rantai pemetaan

\mathdisp {I \stackrel{\partial}{\longrightarrow} I \times \R = TI \stackrel{T(\gamma')}{\longrightarrow} TTM \stackrel{ \pi_{\text{vert} } }{\longrightarrow}  TM} {  }

.



\inputfaktbeweis
{Hyperfläche/Zweite Ableitung einer Kurve/Direkt und über Zusammenhang/Fakt}
{Lema}
{}
{

\faktsituation {Misalkan

\mavergleichskette
{\vergleichskette
{ W
}
{ \subseteq }{ \R^n
}
{  }{
}
{    }{
}
{   }{
}
}
{}{}{}
\definitionsverweis {terbuka}{}{,}
\maabb {h} { W } { \R
} {}
suatu
\definitionsverweis {fungsi yang terdiferensialkan secara kontinu}{}{}
dan

\mavergleichskette
{\vergleichskette
{ Y
}
{ = }{ h^{-1}(c)
}
{  }{
}
{    }{
}
{   }{
}
}
{}{}{}
\definitionsverweis {serat}{}{}
di atas

\mavergleichskette
{\vergleichskette
{ c
}
{ \in }{ \R
}
{  }{
}
{    }{
}
{   }{
}
}
{}{}{,}
dengan $h$ bersifat
\definitionsverweis {reguler}{}{}
di setiap titik $Y$. Misalkan
\maabbdisp {\gamma} { I } { Y
} {}
suatu kurva yang dua kali terdiferensialkan secara kontinu.}
\faktfolgerung {Maka percepatan tangensial $\gamma$ dalam pengertian
Definisi 1.7
sama dengan percepatan tangensial dalam pengertian
Definisi 27.3.}
\faktzusatz {}
\faktzusatz {}

}
{

Menurut
Lema 26.15,
\zusatzklammer {yang juga berlaku dalam dimensi lebih tinggi} {} {},
kedua koneksi tersebut sama.

}

  \zwischenueberschrift{Kurva geodesik}

\inputdefinition
{}
{

Misalkan $M$ suatu
\definitionsverweis {manifold Riemann}{}{}
yang dilengkapi dengan
\definitionsverweis {koneksi Levi--Civita}{}{.}
Suatu
\definitionsverweis {kurva diferensiabel}{}{}
\maabbdisp {\gamma} { I } { M
} {}
yang dua kali terdiferensialkan disebut
\definitionswort {kurva geodesik}{}
\zusatzklammer {atau
\definitionswort {geodesik}{}} {} {,}
apabila

\mavergleichskettedisp
{\vergleichskette
{ \nabla_{d/dt}  \gamma'
}
{ =} { 0
}
{ } {
}
{ } {
}
{ } {
}
}
{}{}{}
berlaku pada $I$.

}

Kadang-kadang digunakan pula istilah \stichwort {geodesik} {.} Perlu ditekankan bahwa kurva geodesik adalah sebuah kurva; dengan demikian, yang dimaksud bukan hanya
\definitionsverweis {citra}{}{}
kurva tersebut, melainkan juga proses geraknya.



\inputfaktbeweis
{Hyperfläche/Geodätische/Direkt und über Zusammenhang/Fakt}
{Lema}
{}
{

\faktsituation {Misalkan

\mavergleichskette
{\vergleichskette
{ W
}
{ \subseteq }{ \R^n
}
{  }{
}
{    }{
}
{   }{
}
}
{}{}{}
\definitionsverweis {terbuka}{}{,}
\maabb {h} { W } { \R
} {}
suatu
\definitionsverweis {fungsi yang terdiferensialkan secara kontinu}{}{}
dan

\mavergleichskette
{\vergleichskette
{ Y
}
{ = }{ h^{-1}(c)
}
{  }{
}
{    }{
}
{   }{
}
}
{}{}{}
\definitionsverweis {serat}{}{}
di atas

\mavergleichskette
{\vergleichskette
{ c
}
{ \in }{ \R
}
{  }{
}
{    }{
}
{   }{
}
}
{}{}{,}
dengan $h$ bersifat
\definitionsverweis {reguler}{}{}
di setiap titik $Y$.}
\faktfolgerung {Maka geodesik-geodesik di $Y$ dalam pengertian
Definisi 1.8
sama dengan geodesik-geodesik dalam pengertian
Definisi 27.5.}
\faktzusatz {}
\faktzusatz {}

}
{

Hal ini mengikuti dari
Lema 27.4.

}



\inputfaktbeweis
{Riemannsche Mannigfaltigkeit/Levi-Civita-Zusammenhang/Geodätische/Geschwindigkeitskonstanz/Fakt}
{Lema}
{}
{

\faktsituation {Misalkan
\maabb {\gamma} { I } { M
} {}
suatu
\definitionsverweis {kurva geodesik}{}{}
pada suatu
\definitionsverweis {manifold Riemann}{}{}
$M$ yang dilengkapi dengan
\definitionsverweis {koneksi Levi--Civita}{}{.}}
\faktfolgerung {Maka
\mathl{\Vert { \gamma'(t) } \Vert}{} konstan pada $I$.}
\faktzusatz {}
\faktzusatz {}

}
{

Pada suatu domain peta terbuka, dengan menggunakan
Teorema 26.14  (2)
dan
Lema 27.2,
diperoleh

\mavergleichskettealign
{\vergleichskettealign
{ { \frac{   \partial }{  \partial t } }  \left\langle  \gamma'(t) ,  \gamma'(t)  \right\rangle
}
{ =} { \left\langle \nabla_{ { \frac{   \partial }{  \partial t } }  }\gamma'(t) ,  \gamma'(t)  \right\rangle + \left\langle  \gamma'(t) , \nabla_{  { \frac{   \partial }{  \partial t } }  }\gamma'(t)  \right\rangle
}
{ =} { 2 \left\langle \nabla_{ { \frac{   \partial }{  \partial t } }  }\gamma'(t) ,  \gamma'(t)  \right\rangle
}
{ =} { 2 \left\langle  0  ,  \gamma'(t)  \right\rangle
}
{ =} { 0
}
}
{}
{}{,}
sehingga
\mathl{\left\langle  \gamma'(t) ,  \gamma'(t)  \right\rangle}{} konstan.

}



\inputfaktbeweis
{Riemannsche Mannigfaltigkeit/Levi-Civita-Zusammenhang/Geodätische/Differentialgleichung/Fakt}
{Teorema}
{}
{

\faktsituation {Suatu
\definitionsverweis {kurva diferensiabel}{}{}
\maabb {\gamma} { I } { M
} {}
yang dua kali terdiferensialkan pada suatu
\definitionsverweis {manifold Riemann}{}{}
$M$ adalah}
\faktfolgerung {suatu
\definitionsverweis {kurva geodesik}{}{}
\zusatzklammer {terhadap
\definitionsverweis {koneksi Levi--Civita}{}{}} {} {,}
jika dan hanya jika pada setiap peta kurva tersebut memenuhi
\definitionsverweis {sistem persamaan diferensial biasa orde kedua}{}{}

\mavergleichskettedisp
{\vergleichskette
{ \gamma_k^{\prime \prime} + \sum_{ij} \Gamma^k_{ij} \gamma_i' \gamma_j'
}
{ =} { 0
}
{ } {
}
{ } {
}
{ } {
}
}
{}{}{}
\zusatzklammer {untuk semua $k$} {} {}.}
\faktzusatz {}
\faktzusatz {}

}
{

Misalkan $n$ dimensi $M$. Kita meninjau situasinya secara langsung pada suatu citra peta terbuka

\mavergleichskette
{\vergleichskette
{ U
}
{ \subseteq }{ \R^n
}
{  }{
}
{    }{
}
{   }{
}
}
{}{}{.}
Menurut
Catatan 25.8,
turunan vertikal diberikan oleh
\maabbeledisp {} {T(U \times \R^n) = U \times \R^n \times \R^n \times \R^n } { U \times \R^n \times  \R^n
} {(P,e_j, \partial_i, w)} {(P,e_j,\sum_{k = 1}^n \Gamma^k_{ij} (P) e_k +w) \cong V(U \times \R^n)
} {},
sedangkan pemetaan ke $TU$ diperoleh dengan menghilangkan komponen tengah. Turunan kedua kurva mula-mula adalah pemetaan tangen kedua, yaitu
\zusatzklammer {dengan mengabaikan perkalian dengan arah berdimensi satu dari ruang tangen kurva} {} {}
\maabbeledisp {T(\gamma)} { I } { TU
} { t } { (\gamma(t), \gamma'(t))
} {,}
dan, sejalan dengan itu,
\maabbeledisp {T(T(\gamma))} { I } { T(TU)
} { t } { ((\gamma(t), \gamma'(t)), (\gamma'(t), \gamma^{\prime \prime} (t)))
} {}
\zusatzklammer {kedua komponen pemetaan tangen diturunkan} {} {.}
Dengan

\mavergleichskette
{\vergleichskette
{ \gamma'(t)
}
{ = }{ \sum_{j = 1}^n \gamma_j'(t) \partial_j
}
{  }{
}
{    }{
}
{   }{
}
}
{}{}{}
hal ini dipetakan oleh proyeksi vertikal ke

\mathdisp {\sum_{k = 1}^n  \left(  \sum_{i,j} \gamma_i'(t) \gamma_j'(t) \Gamma^k_{ij}(\gamma(t))  \right)  \partial_k  +\sum_{k = 1}^n \gamma_k^{\prime \prime}(t) \partial_k} {  }

Syarat bagi suatu geodesik bahwa turunan keduanya
\zusatzklammer {di dalam $TTM$} {} {}
selalu horizontal ekuivalen dengan syarat bahwa proyeksi vertikal yang dihitung di atas sama dengan $0$. Ini berarti setiap komponennya sama dengan $0$, yakni

\mavergleichskettedisp
{\vergleichskette
{ \sum_{i,j} \gamma_i'(t) \gamma_j'(t) \Gamma^k_{ij}(\gamma(t))  + \gamma_k^{\prime \prime}(t)
}
{ =} { 0
}
{ } {
}
{ } {
}
{ } {
}
}
{}{}{}
untuk

\mavergleichskette
{\vergleichskette
{ k
}
{ = }{ 1  , \ldots , n
}
{  }{
}
{    }{
}
{   }{
}
}
{}{}{.}

}

\inputbemerkung
{}
{

Dalam situasi
Teorema 27.8,
sistem persamaan diferensial

\mavergleichskettedisp
{\vergleichskette
{ \gamma_k^{\prime \prime}(t) + \sum_{ij} \Gamma^k_{ij}(\gamma(t)) \gamma_i'(t) \gamma_j'(t)
}
{ =} { 0
}
{ } {
}
{ } {
}
{ } {
}
}
{}{}{}
dengan

\mavergleichskette
{\vergleichskette
{ k
}
{ = }{ 1  , \ldots , n
}
{  }{
}
{    }{
}
{   }{
}
}
{}{}{,}
yang secara lokal mencirikan geodesik, disebut \stichwort {persamaan diferensial geodesik} {.} Ini adalah suatu
\definitionsverweis {sistem persamaan diferensial biasa}{}{}
orde kedua dengan $n$ persamaan. Sistem ini tidak linear dan secara umum sulit diselesaikan.

}

\inputbeispiel{}
{

\bild{ \begin{center}
\includegraphics[width=5.5cm]{ \bildeinlesungsvg {Parallel lines in Poincare's model of hyperbolic geometry} {svg} }
\end{center}
\bildtext {} }

\bildlizenz { Parallel lines in Poincare's model of hyperbolic geometry.svg } {} {Januszkaja} {Commons} {CC-by-sa 3.0} {}

Kita melanjutkan
Contoh 26.7.
Menurut
Teorema 27.8,
suatu geodesik harus memenuhi kedua syarat

\mavergleichskettedisp
{\vergleichskette
{ \gamma_1^{\prime \prime} (t)
}
{ =} { { \frac{   2  }{  \gamma_2(t) } } \gamma_1'(t) \gamma_2'(t)
}
{ } {
}
{ } {
}
{ } {
}
}
{}{}{}
dan

\mavergleichskettedisp
{\vergleichskette
{ \gamma_2^{\prime \prime} (t)
}
{ =} { -  { \frac{   1  }{  \gamma_2(t) } }  \gamma_1'(t) \gamma_1'(t) + { \frac{   1  }{  \gamma_2(t) } }    \gamma_2'(t) \gamma_2'(t)
}
{ } {
}
{ } {
}
{ } {
}
}
{}{}{}.
Untuk

\mavergleichskette
{\vergleichskette
{ \gamma_1
}
{ = }{ x_0
}
{  }{
}
{    }{
}
{   }{
}
}
{}{}{}
yang konstan, sistem ini menjadi

\mavergleichskettedisp
{\vergleichskette
{ \gamma_2^{\prime \prime} (t)
}
{ =} { { \frac{   1  }{  \gamma_2(t) } }  \gamma_2'(t) \gamma_2'(t)
}
{ } {
}
{ } {
}
{ } {
}
}
{}{}{}
dengan solusi komponen

\mavergleichskettedisp
{\vergleichskette
{ y_2(t)
}
{ =} { e^t
}
{ } {
}
{ } {
}
{ } {
}
}
{}{}{}
dan solusi keseluruhan

\mavergleichskettedisp
{\vergleichskette
{ y(t)
}
{ =} { \left(  x_0   , \,  e^t                   \right)
}
{ } {
}
{ } {
}
{ } {
}
}
{}{}{,}
yang bergerak pada suatu garis vertikal. Selain itu, menurut
Soal 27.13,

\mavergleichskettedisp
{\vergleichskette
{ \gamma(t)
}
{ =} { \left(  s   , \,   0                   \right) + r \left(  \tanh  t   , \,   { \frac{   1  }{  \cosh  t  } }                   \right)
}
{ } {
}
{ } {
}
{ } {
}
}
{}{}{}
untuk

\mavergleichskette
{\vergleichskette
{ s
}
{ \in }{ \R
}
{  }{
}
{    }{
}
{   }{
}
}
{}{}{,}

\mavergleichskette
{\vergleichskette
{ r
}
{ \in }{ \R_+
}
{  }{
}
{    }{
}
{   }{
}
}
{}{}{}
merupakan solusi yang, menurut
Soal 20.50 (Analysis (Osnabrück 2021-2023))  (3),
bergerak pada setengah lingkaran dengan pusat
\mathl{\left(  s   , \,   0                   \right)}{} dan jari-jari $r$.

}

\inputbemerkung
{}
{

Apa yang disebut \stichwort {aksioma paralel} {}
\zusatzklammer {atau postulat paralel} {} {}
dalam geometri Euklides bidang menyatakan bahwa untuk setiap garis dan setiap titik yang tidak terletak pada garis tersebut, ada tepat satu garis paralel yang melalui titik itu. Sebuah garis paralel ditentukan oleh sifat bahwa garis tersebut tidak memotong garis yang diberikan. Salah satu pertanyaan penting dalam sejarah adalah apakah postulat paralel dapat diturunkan dari aksioma dan postulat Euklides lainnya, sehingga postulat ini sebenarnya tidak diperlukan. Pada paruh pertama abad ke-19, pertanyaan ini dijawab secara negatif dengan memberikan model-model geometri
\zusatzklammer {yang khususnya memuat konsep titik dan garis} {} {}
yang memenuhi aksioma-aksioma geometri Euklides lainnya, tetapi tidak memenuhi aksioma paralel. Dalam kerangka geometri Riemann, khususnya dengan konsep
\definitionsverweis {geodesik}{}{},
model-model semacam itu muncul secara alami apabila geodesik-geodesik
\zusatzklammer {di sini yang dimaksud adalah citra suatu kurva geodesik} {} {}
dipandang sebagai garis. Perhatikan bahwa dalam pengembangan aksiomatik suatu geometri, garis tidak ditentukan oleh gambaran intuitif, melainkan oleh sifat-sifat yang dimilikinya dalam hubungannya dengan objek-objek lain. Sebagai contoh,
Contoh 27.10
bersama geodesik-geodesik yang dideskripsikan di sana memberikan suatu model geometri non-Euklides: bagi suatu geodesik yang diberikan dan suatu titik lain, selalu terdapat tak terhingga banyak geodesik melalui titik tersebut yang tidak berpotongan dengan geodesik yang diberikan.

}

\inputdefinition
{}
{

Misalkan $M$ suatu
\definitionsverweis {manifold Riemann}{}{}
yang
\definitionsverweis {terhubung}{}{.}
Untuk titik-titik

\mavergleichskette
{\vergleichskette
{ P,Q
}
{ \in }{ M
}
{  }{
}
{    }{
}
{   }{
}
}
{}{}{}
didefinisikan

\mavergleichskettedisp
{\vergleichskette
{ d(P,Q)
}
{ \defeq} { \operatorname{inf} ( L(\gamma), \gamma \text{ adalah lintasan penghubung kontinu yang terdiferensialkan secara kontinu sepotong-sepotong dari } P \text{ ke } Q  )
}
{ } {
}
{ } {
}
{ } {
}
}
{}{}{}
dan besaran ini disebut
\definitionswort {metrik Riemann}{}
pada $M$.

}

Metrik ini benar-benar menjadikan $M$ suatu ruang metrik yang topologinya sama dengan topologi yang diberikan. Selain itu, dapat dibuktikan bahwa realisasi jarak antara dua titik yang diparameterkan oleh panjang busur merupakan suatu geodesik.
